% generated by src.analysis.report_pub from frozen v3 outputs; do not edit by hand
\begin{table*}[tbp]
\centering
\footnotesize
\caption{\textbf{Post hoc descriptive} breakdown of ID-test coverage (\%) by nominal level, at the first and last forecast step (nominal 90\%), and simultaneous whole-trajectory coverage (fraction of windows with all 30 steps and 5 channels inside the interval; no construction targets this quantity). Means over three training repeats. At 90\%, the largest range of coverage across repeats is 0.4 percentage points among the split-conformal rows (Ensemble-normalised conformal, MLP) and 1.2 points among the uncalibrated Gaussian rows (Gaussian raw (no calibration), LSTM); these are seed ranges, not recording-level intervals.}
\label{tab:o1}
\begin{tabular}{lccccccc}
\toprule
Construction & 50\% & 80\% & 90\% & 95\% & Step 1 & Step 30 & Trajectory \\
\midrule
Abs-residual conformal, LSTM & 61.4 & 85.6 & 92.8 & 96.3 & 93.4 & 91.3 & 0.40 \\
Abs-residual conformal, MLP & 61.6 & 85.7 & 92.8 & 96.4 & 93.5 & 91.5 & 0.38 \\
Ensemble-normalised conformal, LSTM & 64.1 & 86.7 & 93.4 & 96.7 & 94.5 & 92.1 & 0.23 \\
Ensemble-normalised conformal, MLP & 64.6 & 86.8 & 93.4 & 96.7 & 94.5 & 91.8 & 0.23 \\
Dropout-normalised conformal, LSTM & 62.5 & 86.7 & 93.5 & 96.7 & 94.3 & 92.2 & 0.36 \\
Dropout-normalised conformal, MLP & 62.7 & 86.3 & 93.2 & 96.6 & 93.7 & 91.9 & 0.31 \\
Sigma-normalised conformal, LSTM & 62.5 & 85.7 & 92.7 & 96.3 & 93.5 & 91.1 & 0.18 \\
Sigma-normalised conformal, MLP & 62.7 & 85.6 & 92.6 & 96.2 & 93.9 & 91.1 & 0.16 \\
Gaussian raw (no calibration), LSTM & 69.8 & 88.9 & 93.7 & 96.0 & 94.7 & 91.6 & 0.22 \\
Gaussian raw (no calibration), MLP & 69.5 & 88.4 & 93.2 & 95.5 & 94.2 & 91.7 & 0.18 \\
\bottomrule
\end{tabular}
\end{table*}
